% -----------------------------------------------------------------------------
% -----------------------------------------------------------------------------

\begin{exbox}[Dyson 自能与寿命]
精确二点函数的谱表示告诉我们，稳定单粒子对应实轴上的孤立极点；一旦该激发能向其它在壳通道衰变，自能获得虚部，极点离开实轴。相应的可计算对象是能量变量 $E$ 下的推迟 Green 函数
\begin{equation*}
 G^R(E)=\frac{1}{E-E_0-\Sigma^R(E)}.
\end{equation*}
在一个窄共振附近，把自能在 $E=E_r$ 邻域展开成
\begin{equation*}
 \Sigma^R(E)
 \simeq \Delta E
 +(1-Z^{-1})(E-E_r)
 -\frac{\ii\Gamma}{2Z},
\end{equation*}
其中 $\Gamma>0$ 具有能量量纲，$Z$ 是极点留数。把实部吸收到重整化共振能量 $E_r=E_0+\Delta E$ 后，Green 函数约化为
\begin{equation}
 \boxed{
 G^R(E)\simeq
 \frac{Z}{E-E_r+\ii\Gamma/2}.}
 \label{eq:dyson-pole-breit-wigner-dp65}
\end{equation}
相应谱函数定义为 $A(E)=-2\operatorname{Im}G^R(E)$，因此
\begin{equation}
 \boxed{
 A(E)=
 \frac{Z\Gamma}{(E-E_r)^2+(\Gamma/2)^2}.}
 \label{eq:dyson-lorentzian-spectral-dp65}
\end{equation}
若该窄共振几乎耗尽单粒子谱权重，则 $Z\simeq1$，并且
\begin{equation*}
 \int_{-\infty}^{\infty}\frac{\dd E}{2\pi}A(E)=Z.
\end{equation*}
所以 $\Gamma$ 不是随意画在谱线上的“展宽参数”，而是由同一个复极点的虚部固定：$\Gamma=-2\,\operatorname{Im}E_\star$，同时也是该 Lorentzian 的半峰全宽。

再把\eqnrefs{eq:dyson-pole-breit-wigner-dp65} Fourier 变回时间。对 $t>0$ 闭合下半平面，只拾取
\begin{equation*}
 E_\star=E_r-\frac{\ii\Gamma}{2}
\end{equation*}
这个极点，得到
\begin{equation*}
 G^R(t)\simeq
 -\frac{\ii Z}{\hbar}\,\theta(t)
 \exp\!\left(-\frac{\ii E_rt}{\hbar}\right)
 \exp\!\left(-\frac{\Gamma t}{2\hbar}\right),
\end{equation*}
因此振幅寿命尺度是 $2\hbar/\Gamma$，而概率或能量按
\begin{equation*}
 \ee^{-\Gamma t/\hbar}
\end{equation*}
衰减，对应通常定义的寿命
\begin{equation}
 \boxed{\tau=\frac{\hbar}{\Gamma}.}
 \label{eq:width-lifetime-si-dp65}
\end{equation}

作为数量级练习，若一个光学共振位于 $E_r=1.55\,\mathrm{eV}$，品质因子 $Q=100$，并按 $Q=E_r/\Gamma$ 定义能量线宽，则
\begin{align*}
 \Gamma&=15.5\,\mathrm{meV},\\
 \tau&=\frac{\hbar}{\Gamma}
 \simeq42.5\,\mathrm{fs}.
\end{align*}
“复极点 $\leftrightarrow$ Lorentzian 频谱 $\leftrightarrow$ 指数时间衰减”是推迟预解式的一般解析结构；有损光学正常模、QNM Green 张量和电子能量损失谱都由相应自能的复极点实现同一结构。


\medskip

稳定粒子对应实轴上的孤立极点；当该态能够衰变时，推迟 Green 函数的极点移入下半复能量平面。\eqnrefs{eq:dyson-pole-breit-wigner-dp65} 已把窄共振附近的分母写成 Breit--Wigner 形式。对能量变量作 Fourier 变换时，下半平面的极点直接产生时间域指数衰减，并给出\eqnrefs{eq:width-lifetime-si-dp65} 的 $\tau=\hbar/\Gamma$。这一近似要求共振孤立、$\Gamma\ll E_r$，且自能在宽度 $\Gamma$ 的能量范围内变化缓慢。


从 Dyson 分母
\begin{equation*}
 D(E)=E-E_0-\Sigma^R(E)
\end{equation*}
开始。把自能实部在 $E_r$ 附近展开：
\begin{equation*}
 \operatorname{Re}\Sigma^R(E)
 \simeq \operatorname{Re}\Sigma^R(E_r)
 +(E-E_r)\,\partial_E\operatorname{Re}\Sigma^R(E_r).
\end{equation*}
令
\begin{equation*}
 E_r-E_0-\operatorname{Re}\Sigma^R(E_r)=0,
 \qquad
 Z^{-1}=1-\partial_E\operatorname{Re}\Sigma^R(E_r),
\end{equation*}
并定义
\begin{equation*}
 \Gamma=-2Z\operatorname{Im}\Sigma^R(E_r)>0.
\end{equation*}
于是
\begin{align*}
 D(E)
 &\simeq Z^{-1}(E-E_r)+\ii\frac{\Gamma}{2Z},\\
 G^R(E)=\frac1{D(E)}
 &\simeq\frac{Z}{E-E_r+\ii\Gamma/2},
\end{align*}
即\eqnrefs{eq:dyson-pole-breit-wigner-dp65}。

取虚部：
\begin{align*}
 \operatorname{Im}\frac{1}{x+\ii a}
 &=-\frac{a}{x^2+a^2},\\
 A(E)=-2\operatorname{Im}G^R(E)
 &=\frac{Z\Gamma}{(E-E_r)^2+(\Gamma/2)^2},
\end{align*}
得到\eqnrefs{eq:dyson-lorentzian-spectral-dp65}。再计算时间核
\begin{equation*}
 G^R(t)=\int_{-\infty}^{\infty}\frac{\dd E}{2\pi\hbar}
 \ee^{-\ii Et/\hbar}G^R(E).
\end{equation*}
对 $t>0$，因 $\ee^{-\ii Et/\hbar}$ 在下半平面衰减，闭合下半圆；唯一极点是 $E_\star=E_r-\ii\Gamma/2$。留数定理给
\begin{equation*}
 G^R(t)\simeq
 -\frac{\ii Z}{\hbar}\,
 \ee^{-\ii E_rt/\hbar}\ee^{-\Gamma t/(2\hbar)}.
\end{equation*}
振幅模方因而为
\begin{equation*}
 |G^R(t)|^2\simeq
 \frac{|Z|^2}{\hbar^2}\ee^{-\Gamma t/\hbar}.
\end{equation*}
若把寿命定义为概率下降到 $\ee^{-1}$ 的时间，便得到\eqnrefs{eq:width-lifetime-si-dp65}。
\end{exbox}


